\documentclass[11pt,reqno]{amsart}
\usepackage[T1]{fontenc}
\usepackage{amsmath,amssymb,mathtools}
\usepackage[margin=1in]{geometry}
\usepackage[colorlinks=true,linkcolor=blue,citecolor=blue,urlcolor=blue]{hyperref}
\setlength{\emergencystretch}{2em}
\title{Induced Multinomial Chains, Cumulative Trials, and the Mehler Limit}
\date{Updated sections from the Kravchuk Encyclopedia, Chapters 3 and 11}
\begin{document}\maketitle
\section{Induced matrices and multinomial chains: the SFA II framework}
\label{sec:ency-sfa-ii}
The multinomial extension of a one-site chain is the induced-matrix
construction of SFA II \cite{FeinsilverSFAII2023}. It provides the finite
framework for the cumulative-trial example below and the Mehler limit
in Section~\ref{sec:11-3}. We declare row-stochastic conventions here;
column-vector probability evolution uses the transpose.

For a $q$-state matrix $A$, its degree-$N$ induced matrix is defined by
\begin{equation}\label{eq:ency-induced-definition}
\prod_a\left(\sum_bA_{ab}z_b\right)^{n_a}
=\sum_{|\mathbf m|=N}A^{[N]}(\mathbf n,\mathbf m)\mathbf z^{\mathbf m},
\qquad |\mathbf n|=N.
\end{equation}
Substitution gives $(AB)^{[N]}=A^{[N]}B^{[N]}$. If $A=P$ is stochastic,
$P^{[N]}$ is the transition for $N$ independent sites with their labels
forgotten. Equivalently,
\[
P^{[N]}(\mathbf n,\mathbf m)
=\sum_{\substack{\sum_b c_{ab}=n_a\\\sum_a c_{ab}=m_b}}
\prod_a\frac{n_a!}{\prod_b c_{ab}!}\prod_{a,b}P_{ab}^{c_{ab}}.
\]
Thus the contingency-table model and the induced matrix are exactly
the same finite construction. The coordinate matrix of the symmetric
tensor power depends on basis normalization; the stochastic occupation
basis should not be confused with an orthonormal tensor basis.

If $\boldsymbol\eta$ is stationary, then
\[
\Pi_N(\mathbf n)=\frac{N!}{\prod_a n_a!}\prod_a\eta_a^{n_a}
\]
is stationary. For a complete one-site eigenbasis with eigenvalues
$1,r_1,\ldots,r_{q-1}$, symmetric products give occupation eigenvalues
$\prod_{j=1}^{q-1}r_j^{\nu_j}$ for $|\boldsymbol\nu|\le N$.
Reversibility permits orthogonal one-site modes and the associated
multivariate Krawtchouk eigenbasis; stochasticity alone does not imply
this orthogonality.

For $P_t=e^{tQ}$ the induced generator is
\[
(\mathcal L_Nf)(\mathbf n)=\sum_{a\ne b}n_aQ_{ab}
[f(\mathbf n-e_a+e_b)-f(\mathbf n)].
\]
Cumulative Bernoulli trials specify a particular one-site $P$ and
therefore a particular $P^{[N]}$; they are not a different multinomial
extension. Their sampling interpretation makes positivity transparent,
while the induced construction explains composition and spectrum.
Matrix composition, polynomial linearization, and positive radial
convolution are distinct products, related only through their specified
spectral realizations.

Under stationary input the joint site probabilities are
$J_{ab}=\eta_aP_{ab}$. Recording $N$ such pairs gives a multinomial table
with margins $\mathbf N,\mathbf M$, and
\[
\frac{P^{[N]}(\mathbf n,\mathbf m)}{\Pi_N(\mathbf m)}
=\frac{\Pr(\mathbf N=\mathbf n,\mathbf M=\mathbf m)}
{\Pi_N(\mathbf n)\Pi_N(\mathbf m)}.
\]
This density ratio is the finite kernel whose central local limit is
Mehler's kernel. In the reversible case, diagonal fluctuation modes
give $q-1$ Gaussian factors and, after conjugation, $q-1$ oscillator
number operators. These limiting assertions require the local-limit
and mode-contraction arguments, beyond the induced-matrix identity.

\section{Cumulative Bernoulli trials: saved and recouped successes}
\label{sec:cumulative-bernoulli}
Hoare and Rahman \cite{HoareRahman1984} develop Krawtchouk processes
from combinations of Bernoulli trials with saved and recouped successes.
The earlier dynamic-occupancy program of Cooper, Hoare, and Rahman
\cite{CooperHoareRahman1977} starts from a Hahn kernel with
negative-hypergeometric factors. The following independent-site
construction is the binary specialization of Section~\ref{sec:ency-sfa-ii}
and presents our Bernoulli case with explicit conventions;
its parameters are defined here rather than implicitly identified with
all conventions in those papers.

Starting from $i$ successes among $N$ sites, save each success with
probability $\alpha$. If $k$ are saved, retry the remaining $N-k$ sites
with success probability $\beta$. With
$b(j;m,p)=\binom mjp^j(1-p)^{m-j}$ (zero outside $0\le j\le m$), the
positive transition kernel is
\begin{equation}\label{eq:ency-cumulative}
\mathcal K(j\mid i)=\sum_{k=0}^{\min(i,j)}
b(k;i,\alpha)b(j-k;N-k,\beta).
\end{equation}
Its backward action is $(\mathcal Kf)(i)=\sum_j\mathcal K(j\mid i)f(j)$;
the column-vector forward transition matrix is its transpose.
For $0\le\alpha<1$, $0<\beta<1$, the one-site transition and its
stationary parameter are
\[
P=\begin{pmatrix}1-\beta&\beta\\
(1-\alpha)(1-\beta)&\alpha+(1-\alpha)\beta\end{pmatrix},\qquad
r=\alpha(1-\beta),\quad p=\frac{\beta}{1-r}.
\]
Indeed $\mathbb E[Y-p\mid X]=r(X-p)$. The invariant occupation law is
$w_{N,p}(i)=\binom Ni p^i(1-p)^{N-i}$, and the biased Krawtchouk
coefficients defined by
\[
(1+(1-p)z)^i(1-pz)^{N-i}=\sum_{n=0}^NC_n(i;N,p)z^n
\]
satisfy
\[
\mathcal KC_n=r^nC_n,\qquad
\sum_iw_{N,p}(i)C_n(i)C_m(i)=\binom Nn[p(1-p)]^n\delta_{nm}.
\]
Both statements follow by expanding products of centered site variables
$X_a-p$. Thus the positive formula \eqref{eq:ency-cumulative} has the
exact bilinear expansion
\begin{equation}\label{eq:ency-cumulative-spectral}
\frac{\mathcal K(j\mid i)}{w_{N,p}(j)}
=\sum_{n=0}^N\frac{r^nC_n(i;N,p)C_n(j;N,p)}{\binom Nn[p(1-p)]^n}.
\end{equation}

For our symmetric convention, take $p=1/2$ and
$\beta=(1-r)/2$, $\alpha=2r/(1+r)$, with $0\le r<1$.
Then $P=\frac12\left(\begin{smallmatrix}1+r&1-r\\1-r&1+r\end{smallmatrix}\right)$
and $C_n(i;N,1/2)=(-1/2)^nK_n(i;N)$ in ordinary coefficient
normalization. For $r=e^{-t}$ this is the time-$t$ independent-bit
transition. It is not the single-flip Ehrenfest step: the spectra are
$r^n$ and $1-2n/N$, respectively.

At stationarity the $N$ input--output site pairs give a four-category
multinomial table. Its input and output margins $I,J$ obey
\[
\frac{\mathcal K(j\mid i)}{w_{N,p}(j)}
=\frac{\Pr(I=i,J=j)}{w_{N,p}(i)w_{N,p}(j)}.
\]
This is the symmetric-tensor occupation kernel, interpreted as a joint
law relative to the product of its margins. Conditioning instead on a
fixed flip count gives hypergeometric overlap laws. Central Gaussian
scaling of the joint density ratio produces the Mehler kernel in
Section~\ref{sec:11-3}.

\section{Bilinear kernels and the Mehler limit}
\label{sec:11-3}
The finite foundation is the induced multinomial chain of SFA II
\cite{FeinsilverSFAII2023}, presented in
Section~\ref{sec:ency-sfa-ii}. The Gaussian and oscillator constructions
below continue that framework in a central fluctuation limit; the
cumulative-trial model is a particular stochastic realization.

The saved-success model of Section~\ref{sec:cumulative-bernoulli}
gives a finite positive-kernel derivation of the Mehler limit,
following the dynamic-occupancy viewpoint
\cite{CooperHoareRahman1977,HoareRahman1984}. Fix $0<r<1$ in its
symmetric case. Write
\[
w_N(k)=2^{-N}\binom Nk,\quad
p_{n,N}(k)=\binom Nn^{-1/2}K_n(k;N),\quad
(1+z)^{N-k}(1-z)^k=\sum_nK_n(k;N)z^n.
\]
The finite kernel relative to $w_N$ is
\begin{equation}\label{eq:ency-finite-mehler}
m_{r,N}(k,\ell)=\sum_{n=0}^Nr^np_{n,N}(k)p_{n,N}(\ell)
=\sum_{n=0}^N\frac{r^nK_n(k;N)K_n(\ell;N)}{\binom Nn}.
\end{equation}
For stationary input and output counts $K,L$,
$m_{r,N}(k,\ell)=\Pr(K=k,L=\ell)/(w_N(k)w_N(\ell))$ and
\[
\mathbb E[u^Kv^L]=
\left[\tfrac{1+r}{4}(1+uv)+\tfrac{1-r}{4}(u+v)\right]^N.
\]
Set $x_k=(N-2k)/\sqrt{2N}$ and $\delta_N=\sqrt{2/N}$.
The joint local central limit theorem gives, on compact central windows,
\[
\frac{\Pr(K=k,L=\ell)}{\delta_N^2}\longrightarrow
\frac{\exp[-(x^2-2rxy+y^2)/(1-r^2)]}{\pi\sqrt{1-r^2}},
\qquad \frac{w_N(k)}{\delta_N}\longrightarrow
w(x)=\pi^{-1/2}e^{-x^2}.
\]
Consequently
\begin{equation}\label{eq:ency-mehler-ratio}
m_{r,N}(k,\ell)\longrightarrow m_r(x,y)
=\frac1{\sqrt{1-r^2}}
\exp\left[\frac{2rxy-r^2(x^2+y^2)}{1-r^2}\right].
\end{equation}
The density is relative to $d\mu=w(x)dx$, of variance $1/2$.
The transition is $Y=rx+\sqrt{1-r^2}Z$ for $Z\sim\mu$.
For $r=e^{-t}$ its generator is $G=\frac12D^2-xD$.
The unitary $Uf=w^{1/2}f$ conjugates $-G$ to
$H-\frac12I$, where $H=\frac12(x^2-D^2)$. Thus
\[
\sqrt{w(x)}m_r(x,y)\sqrt{w(y)}
=\frac{\exp[-(1+r^2)(x^2+y^2)/(2(1-r^2))+2rxy/(1-r^2)]}
{\sqrt{\pi(1-r^2)}}
\]
is the kernel of $r^{H-1/2}$, and multiplication by $\sqrt r$ gives
$r^H$. This is Mehler's kernel in oscillator coordinates.

In the probabilists' Hermite convention,
the bilinear identity is
\[
\sum_{n\ge0}\frac{r^n}{n!}\operatorname{He}_n(\xi)
\operatorname{He}_n(\eta)
=\frac1{\sqrt{1-r^2}}
\exp\left[\frac{2r\xi\eta-r^2(\xi^2+\eta^2)}{2(1-r^2)}\right],
\qquad |r|<1.
\]
It converges absolutely in that disk, with the square root equal to
$1$ at $r=0$. The fixed-degree Krawtchouk--Hermite limit identifies the
spectral modes; the local-limit argument above proves convergence of
the finite kernels without an unproved interchange of an expanding
spectral sum with a limit. Neither assertion alone proves convergence
of an unrelated weighted finite-matrix spectrum.

\begin{thebibliography}{9}
\bibitem{FeinsilverSFAII2023}
P.~Feinsilver, \emph{Symmetric Functions Algebras (SFA) II: Induced
Matrices}, Journal of Stochastic Analysis \textbf{4} (2023), no.~2,
Article~1.
\href{https://doi.org/10.31390/josa.4.2.01}{doi:10.31390/josa.4.2.01}.

\bibitem{CooperHoareRahman1977}
R.~D. Cooper, M.~R. Hoare, and M.~Rahman,
\emph{Stochastic processes and special functions: On the probabilistic
origin of some positive kernels associated with classical orthogonal
polynomials}, J. Math. Anal. Appl. \textbf{61} (1977), no.~1, 262--291.
\href{https://doi.org/10.1016/0022-247X(77)90160-3}{doi:10.1016/0022-247X(77)90160-3}.

\bibitem{HoareRahman1984}
M.~R. Hoare and M.~Rahman, \emph{Cumulative Bernoulli trials and
Krawtchouk processes}, Stochastic Process. Appl. \textbf{16} (1984),
no.~2, 113--139.
\href{https://doi.org/10.1016/0304-4149(84)90014-0}{doi:10.1016/0304-4149(84)90014-0}.
The publisher dates the issue February 1984; its copyright notice is 1983.

\end{thebibliography}\end{document}